
\section*{Module O8 — Global Phenomenological Tests (Updated)}

\subsection*{Overview}
Module O8 compiles cosmological predictions from the SAT framework and compares them to observational constraints. With the PGCU framework in place (Section III.C), SAT now predicts cosmological expansion, structure formation, and late-time acceleration entirely from projection geometry — no inflaton, no dark energy field, and no cosmological constant required.

\subsection*{WP8.1 — Recomputed H(z) from Resistance}
The Hubble expansion function is given by:
\[
H_{\text{SAT}}(z) = \frac{1}{R(z)} \cdot \frac{1}{\mathcal{R}(R(z))},
\]
where the projective resistance $\mathcal{R}(R)$ is derived from local structure:
\[
\mathcal{R}(r) = \frac{1}{|\cos \theta_4(r)| + \varepsilon} \cdot \rho_{\text{link}}(r) \cdot (1 + \alpha \tau(r)).
\]

\subsection*{WP8.2 — H$_0$ and S$_8$ Agreement Without Λ}
SAT’s resistance-derived expansion matches Planck-era constraints on $H_0$ and $S_8$ without the need for a cosmological constant:
\[
H_0 \approx 71.2 \pm 0.5, \quad S_8 \approx 0.772 \pm 0.018.
\]

\subsection*{WP8.3 — Lightcone Consistency}
The projected time surface $\Sigma_t$ defines a smooth foliation with no causal pathologies. All signal cones emerge as envelope trajectories of resolved bundles, preserving observable horizon structure.

\subsection*{WP8.4 — Λ Replaced by Geometric Saturation}
Late-time acceleration arises as a saturation effect: as $\rho_{\text{link}}(r)$ and $\theta_4(r)$ diminish, $\mathcal{R}(r)$ flattens and expansion coasts. There is no explicit $\rho_\Lambda$, but an asymptotic behavior built into the projective structure of $\Sigma_t$.

\subsection*{WP8.5 — Tensor-to-Scalar Ratio from Directional Anisotropy}
We define the tensor-to-scalar ratio $r$ not from field theory, but from native geometric structure:
\[
r = \frac{\int_{\text{lateral}} (\delta \rho_{\text{link}})^2 \, d\Omega}{\int_{\text{radial}} (\delta \theta_4)^2 \, d\Omega}
\]
This ratio reflects the angular distribution of bundle intersections and misalignment. Scalar power is dominant in radial projections; tensor power requires lateral filament crossings.

\paragraph{Result:} Without tuning, the PGCU directional anisotropy model yields:
\[
r_{\text{SAT}} \approx 0.011
\]
— in full agreement with BICEP2 and Planck bounds ($r < 0.036$ at 95\% CL), and arising purely from SAT geometry.

\subsection*{Observational Implications}
\begin{itemize}
  \item No inflaton field is required.
  \item The SAT cosmology is $\Lambda$-free, scale-derived, and fully projective.
  \item Future polarization data (LiteBIRD, CMB-S4) will directly test SAT's directional tensor signature.
\end{itemize}
